\section{尾声：贡献、限制与结论}

% ---------------------------------------------------------------- 两点贡献
\begin{frame}{两点贡献}
  \begin{block}{贡献 (1)：揭示并刻画零空间漂移 + 输入侧补救}
    \small
    合成训练的网络迁移至真实数据后，诊断证据显示其新增高频与观测算子的\textbf{零/近零空间}漂移一致；
    结合固定算子的零空间分析，\redflag{任何有界观测侧锚定项对该漂移分量均不可见}——
    有效补救不在更强惩罚，而在\textbf{输入侧 drizzle 证据注入}。
  \end{block}
  \begin{block}{贡献 (2)：无真值协议下经典 vs 学习的诚实裁决}
    \small
    以相位分层 split-half FRC（带控制组）为相干信息定界，辅以一对观测锚定 proxy 与机械化选点；
    该判据\textbf{本身是有界的 proxy，而非光学真值}。在此判据下：经典变分重建最忠于可验证观测，
    学习重建表观轮廓最丰富但携带最多不可验证高频——\redflag{不存在 GT-可认证的单一最优方法}。
  \end{block}
  \takeaway{贡献 (1) 是机制发现，贡献 (2) 是在\textbf{诚实判据}下的取舍，而非宣布赢家。}
\end{frame}

% ---------------------------------------------------------------- 限制
\begin{frame}{限制：把不能声称的话也说清楚}
  \begin{enumerate}
    \small
    \item \textbf{单仪器、单场景族}：全部真实数据来自一台标定系统上的一颗芯片（248 帧）；
          定位为\textbf{标定化 case study + 可复用协议}，外部效度是论证而非演示。
    \item \textbf{无外部高分辨率锚}：没有独立的更高分辨率热成像验证；
          “可信”指\textbf{内部门控下的一致性}，而非与持有真值一致。
    \item \textbf{FRC 控制组部分失常}：$10$--$12\,$µm 回弹周围正/负控制行为不理想；
          缓解是只主张 $17\,$µm cutoff 并把回弹标为伪影风险。
    \item \textbf{Proxy 天花板}：proxy 能为幻觉设防，却\textbf{不能认证}恢复出的细节；
          采纳规则保守，可能拒绝真正有用的重建。
    \item \textbf{最细结构的物理上限}：$10\,$µm pitch 下 $1$--$2\,$px 线位于分辨率极限，
          叠加 PSF 不确定度（$\sigma\in[0.2,0.5]$），交付物是\textbf{可见性与稳定性，而非复原}。
  \end{enumerate}
  \takeaway{未来工作：\textbf{多场景采集协议}已有设计；引入\textbf{外部热成像锚}的验证价值最高。}
\end{frame}
